\scanpage{47}
non plus le théorème de dualité, ou de Künneth. \deleted{Je \unsure{ne} \unsure{d'autre part}}
(Il semble cependant que malgré cela, ils arrivent à utiliser leur cohomologie pour donner une expression à la Lefschetz-Weil de la fonction $\zeta$ d'un schéma projectif et lisse sur un corps fini, en montrant que les Frobenius opérant sur les $H^i(X)$ sont des opérateurs à trace .... ! Le formalisme semble alors proche de celui de Dwork, que je ne connais pas .... ).

\S\ \unsure{3} \inserted{\unsure{4)}} \uline{Relations \inserted{\uline{heuristique de} W.M.} avec mes idées de Gauss-Manin.}

1. Caractère assez surprenant à première vue de l'assertion d'invariance de W.M. \unsure{Supposons} p.ex. que l'on ait

\margin{46}
